% =====================================================================
%  1. Why gradient descent?
% =====================================================================
\section{Why gradient descent?}

\begin{frame}{The question behind it: fit a line to data}
  \begin{columns}[T]
    \begin{column}{0.5\textwidth}
      \small
      Predict height from weight with a straight line, $\hat y = b + m x$
      (intercept $b$, slope $m$).
      How good is a line? Add up the squared \alert{residuals} $r_i = y_i - (b + m x_i)$:
      \[
        \SSR(b, m) = \sum_{i=1}^{n}\bigl(y_i - b - m x_i\bigr)^2 .
      \]
      This is the \textbf{loss function}. The best line is the one with the smallest SSR.
      \begin{keybox}[Other uses]
        \footnotesize Linear and logistic regression, t-SNE: each one is ``pick the parameters that make a loss smallest''.
      \end{keybox}
    \end{column}
    \begin{column}{0.48\textwidth}
      \centering
      \only<1>{\includegraphics[width=\linewidth,height=0.55\textheight,keepaspectratio]{fig_line_frame_0}}%
      \only<2>{\includegraphics[width=\linewidth,height=0.55\textheight,keepaspectratio]{fig_line_frame_1}}%
      \only<3>{\includegraphics[width=\linewidth,height=0.55\textheight,keepaspectratio]{fig_line_frame_2}}%
      \only<4>{\includegraphics[width=\linewidth,height=0.55\textheight,keepaspectratio]{fig_line_frame_3}}%
      \only<5>{\includegraphics[width=\linewidth,height=0.55\textheight,keepaspectratio]{fig_line_frame_4}}%
      \only<6>{\includegraphics[width=\linewidth,height=0.55\textheight,keepaspectratio]{fig_line_frame_5}}%

      \footnotesize Data: weights $0.5, 2.3, 2.9$, heights $1.4, 1.9, 3.2$. Each click shows smaller residuals.
    \end{column}
  \end{columns}
\end{frame}

\begin{frame}{Why not just set the slope of the loss to zero?}
  \begin{columns}[T]
    \begin{column}{0.48\textwidth}
      \small
      At the best value the loss is flat, so solve loss$'=0$. For a line we \emph{can} (least squares).
      But often we \alert{cannot}: there is no closed form (logistic regression, neural nets),
      there are millions of parameters, or new data keeps arriving.
      \begin{thmbox}[The idea of gradient descent]
        \footnotesize Start from a guess and step downhill until the loss stops improving:
        \alert{big steps} far away, \alert{small steps} near the optimum.
        Cauchy proposed it in 1847~\cite{cauchy1847}; it is still used to train large models~\cite{bottou2018}.
      \end{thmbox}
    \end{column}
    \begin{column}{0.5\textwidth}
      \centering
      \includegraphics[width=\linewidth, height=0.55\textheight, keepaspectratio]{fig_ssr_intercept}

      \footnotesize SSR against the intercept (slope $0.64$). Steps shrink as the slope flattens.
    \end{column}
  \end{columns}
\end{frame}
